% creation date: may 11, 2026
% last modified: may 11, 2026
% description: latex source for linked list cycle detection summary
% author: enigmak9

\documentclass{article}
\usepackage[utf8]{inputenc}
\usepackage{amsmath}
\usepackage{geometry}

\geometry{margin=1in}

\begin{document}

\title{Linked List Cycle Detection Summary}
\author{}
\date{May 11, 2026}
\maketitle

\section*{Overview}
% summary of cycle detection logic
This document outlines the techniques used to identify cycles within a linked list, specifically comparing the hash set method with the more space efficient tortoise and hare algorithm.

\section*{1. Hash Set Approach}
% explanation of the visit tracking method
The most straightforward method for cycle detection is to track every node visited during traversal.
\begin{itemize}
    \item \textbf{Mechanism}: Use a hash set to store node references.
    \item \textbf{Detection}: If a node being visited is already present in the set, a cycle is confirmed.
    \item \textbf{Time Complexity}: $O(n)$
    \item \textbf{Space Complexity}: $O(n)$
\end{itemize}

\section*{2. Floyd's Tortoise and Hare Algorithm}
% explanation of the two pointer technique
This algorithm optimizes space by using two pointers that move at different speeds.
\begin{itemize}
    \item \textbf{Slow Pointer}: Moves 1 step at a time.
    \item \textbf{Fast Pointer}: Moves 2 steps at a time.
\end{itemize}
If the fast pointer reaches the end of the list (null), there is no cycle. If it eventually meets the slow pointer at the same node, a cycle exists.

\section*{Mathematical Convergence}
% proof of linear time complexity
The convergence of the pointers is guaranteed by the decreasing gap between them. Let $G$ be the initial integer distance between the pointers within a cycle. In each iteration:
\begin{enumerate}
    \item The slow pointer moves forward by 1 unit.
    \item The fast pointer moves forward by 2 units.
\end{enumerate}

The resulting gap for the next iteration is:
\[ G_{new} = G + 1 - 2 = G - 1 \]

Since the gap decreases by exactly 1 in each step, it will reach 0 in at most $n$ steps, where $n$ is the cycle length. This ensures the algorithm runs in linear time $O(n)$ with constant space $O(1)$.

\end{document}
